\documentclass[11pt,a4paper]{article}

% ---------------------------------------------------------------------------
% Packages
% ---------------------------------------------------------------------------
\usepackage{amsmath,amssymb}
\usepackage{physics}
\usepackage{braket}
\usepackage{hyperref}
\usepackage{graphicx}
\usepackage{booktabs}
\usepackage{geometry}
\geometry{margin=2.5cm}

% ---------------------------------------------------------------------------
% Custom commands for this note
% ---------------------------------------------------------------------------
\newcommand{\eps}{\varepsilon}
\newcommand{\Gint}{\Gamma_{\text{interp}}}
\newcommand{\Gsink}{\Gamma_{\text{sink}}}
\newcommand{\PolProj}{\mathcal{P}}
\newcommand{\Cmat}{C}
\newcommand{\gfive}{\gamma_{5}}
\newcommand{\gp}{\gamma_{+}}
\newcommand{\gm}{\gamma_{-}}
\newcommand{\Dmu}{D_{\mu}}
\newcommand{\Dnum}{D_{\nu}}
\newcommand{\Tmunu}{T_{\mu\nu}}
\newcommand{\TmunuQ}{T_{\mu\nu}^{q}}
\newcommand{\TmunuG}{T_{\mu\nu}^{g}}
\newcommand{\seqsrc}{\eta_{\text{seq}}}
\newcommand{\seqprop}{\Sigma}
\newcommand{\fwprop}{S}
\newcommand{\bwprop}{\bar{S}}
\newcommand{\cf}{c_{f}}
\newcommand{\cs}{c_{s}}
\newcommand{\xsrc}{x_{0}}
\newcommand{\xsink}{y}
\newcommand{\xins}{x}
\newcommand{\pfinal}{p_{f}}
\newcommand{\pini}{p_{i}}
\newcommand{\momtrans}{q}
\newcommand{\code}[1]{\texttt{#1}}

% ---------------------------------------------------------------------------
\title{\textbf{Proton Energy-Momentum Tensor Measurement}\\[4pt]
       \large A Detailed Physics Note for the PyQUDA Implementation}
\author{}
\date{\today}

\begin{document}
\maketitle

\begin{abstract}
This note documents the complete set of formulas, conventions, and algorithmic steps
used in the PyQUDA-based measurement of the proton matrix elements of the
energy-momentum tensor (EMT).  We cover the proton interpolating field, the connected
two-point function, the fixed-sink sequential-source method for connected three-point
functions, the covariant derivative insertion for the quark EMT, gradient flow, and the
output data structures.  The note is written for lattice QCD practitioners and aims to
be sufficiently precise that every einsum in the code can be traced back to a
well-defined mathematical expression.  Throughout, we adopt the Euclidean metric with
Hermitian gamma matrices and the DeGrand--Rossi chiral basis as implemented in QUDA
and exposed through PyQUDA.
\end{abstract}

\tableofcontents

% ===========================================================================
\section{Conventions and Gamma Matrix Setup}
% ===========================================================================

\subsection{Euclidean Gamma Matrices}

We work in four-dimensional Euclidean spacetime with coordinates
$x_{\mu}$, $\mu = 1,2,3,4$.  The spatial directions are $1,2,3$ and the temporal
direction is $\mu = 4$.  The Euclidean gamma matrices satisfy
\begin{equation}
  \{\gamma_{\mu}, \gamma_{\nu}\} = 2\delta_{\mu\nu}\,\mathbb{1}_{4\times 4},
  \qquad \gamma_{\mu}^{\dagger} = \gamma_{\mu}.
\end{equation}
The chiral matrix is
\begin{equation}
  \gamma_{5} \equiv \gamma_{1}\gamma_{2}\gamma_{3}\gamma_{4},
  \qquad \gamma_{5}^{\dagger} = \gamma_{5},\quad \gamma_{5}^{2} = \mathbb{1}.
\end{equation}

\subsection{PyQUDA Gamma Indexing}
\label{sec:gammas}

PyQUDA uses the QUDA bitmask convention for gamma matrices.  The 16 independent
bilinear structures are ordered as
\begin{equation}
  \texttt{my\_gammas} = [\texttt{"5"}, \texttt{"T"}, \texttt{"T5"}, \texttt{"X"},
  \texttt{"X5"}, \texttt{"Y"}, \texttt{"Y5"}, \texttt{"Z"}, \texttt{"Z5"},
  \texttt{"I"}, \texttt{"SXT"}, \texttt{"SXY"}, \texttt{"SXZ"}, \texttt{"SYT"},
  \texttt{"SYZ"}, \texttt{"SZT"}].
\end{equation}
The corresponding QUDA-internal integer identifiers are
\begin{equation}
  \texttt{pyquda\_gammas\_order} = [15,\,8,\,7,\,1,\,14,\,2,\,13,\,4,\,11,\,
  0,\,9,\,3,\,5,\,10,\,6,\,12].
\end{equation}
The Dirac gamma matrices relevant for derivative insertions are
\begin{equation}
  \gamma_{1}\leftrightarrow\texttt{gamma(1)},\quad
  \gamma_{2}\leftrightarrow\texttt{gamma(2)},\quad
  \gamma_{3}\leftrightarrow\texttt{gamma(4)},\quad
  \gamma_{4}\leftrightarrow\texttt{gamma(8)}.
\end{equation}
The chiral matrix is $\gamma_{5}\leftrightarrow\texttt{gamma(15)}$.

\subsection{Charge Conjugation Matrix}

The charge conjugation matrix in the Euclidean DeGrand--Rossi basis is
\begin{equation}
  \Cmat = i\gamma_{2}\gamma_{4}
        = i\;\texttt{gamma(2)}\;\texttt{ @ }\;\texttt{gamma(8)}.
\end{equation}
It satisfies the standard relations
\begin{equation}
  \Cmat\gamma_{\mu}^{T}\Cmat^{-1} = -\gamma_{\mu},\qquad
  \Cmat^{T} = -\Cmat,\qquad
  \Cmat^{\dagger} = \Cmat^{-1} = -\Cmat.
\end{equation}

\subsection{Propagator Data Layout}

A lattice quark propagator in PyQUDA is stored with the index structure
\begin{equation}
  \fwprop(x,x_{0}) \;\longleftrightarrow\;
  \texttt{[w,\, t,\, z,\, y,\, x\_cb,\, spin\_sink,\, spin\_src,\, color\_sink,\, color\_src]},
\end{equation}
where \texttt{w} is the checkerboard/parity index, (\texttt{t,z,y,x\_cb}) are the
local spacetime coordinates, and the remaining four indices carry the spin and color
degrees of freedom of the two ends of the propagator.  In einsum notation we label
this layout as
\begin{equation}
  \fwprop^{i,k}_{a,d}(x,x_{0}) \quad\text{with}\quad
  i = \text{spin\_sink},\; k = \text{spin\_src},\;
  a = \text{color\_sink},\; d = \text{color\_src}.
\end{equation}
Greek letters from the middle of the alphabet ($\alpha,\beta,\gamma,\dots$) denote
spin indices and Latin letters $(a,b,c,\dots)$ denote color indices when we write
mathematical formulas.  In the einsum code listings, spin indices are
$i,j,k,l,m,n$ and color indices are $a,b,c,d,e,f$.

\subsection{Momentum Phases}

Momentum projection uses the \texttt{MomentumPhase} class from PyQUDA.  For a
momentum $k = (k_{1},k_{2},k_{3},0)$ and source position $\xsrc$, the phase at
position $r$ is
\begin{equation}
  \Phi_{k}(r-\xsrc) = \exp\!\big(-i\,k\cdot(r-\xsrc)\big),
\end{equation}
where the sign convention follows the \texttt{MomentumPhase} implementation and
the dot product uses the spatial components only.

% ===========================================================================
\part{Proton Interpolating Field and Two-Point Function}
% ===========================================================================

% ---------------------------------------------------------------------------
\section{Proton Interpolating Field}
% ---------------------------------------------------------------------------

The proton is interpolated by a color-singlet three-quark operator with
spin-$\tfrac{1}{2}$.  The standard choice used in the code is the scalar-diquark
interpolator.  Expressed in terms of up and down quark fields, the proton
interpolating field at position $y$ is
\begin{equation}
  \boxed{
  \chi_{\alpha}(y) = \eps_{abc}\;
  \big[u_{a}^{T}(y)\,\Cmat\,\Gint\,d_{b}(y)\big]\;
  u_{c,\alpha}(y)
  },
  \label{eq:chi}
\end{equation}
where:
\begin{itemize}
  \item $\alpha = 1,\dots,4$ is the free Dirac index of the proton;
  \item $a,b,c = 1,2,3$ are fundamental color indices;
  \item $\eps_{abc}$ is the totally antisymmetric Levi-Civita tensor with
        $\eps_{123} = +1$;
  \item $\Cmat = i\gamma_{2}\gamma_{4}$ is the charge conjugation matrix;
  \item $\Gint$ selects the spin structure of the diquark channel;
  \item $u(x), d(x)$ are the up- and down-quark Dirac spinor fields.
\end{itemize}

The transpose $u_{a}^{T}$ acts in Dirac space only; the color index $a$ is
contracted via the epsilon tensor.  The bracketed term $[u^{T}\Cmat\Gint d]$
forms a Lorentz scalar (for $\Gint = \gamma_{5}$) or a pseudoscalar diquark, and
the remaining up quark $u_{c,\alpha}$ carries the spin of the proton.

\subsection{Supported Interpolators}

The code supports three choices for $\Gint$, mapped by the string argument
\texttt{interpolator}:

\begin{center}
\begin{tabular}{ccl}
\toprule
\texttt{interpolator} & $\Gint$ & Description \\
\midrule
\texttt{"5"}  & $\Cmat\gamma_{5}$   & Scalar diquark (standard) \\
\texttt{"T5"} & $\Cmat\gamma_{t}\gamma_{5}$ & Temporal-axial diquark \\
\texttt{"Z5"} & $\Cmat\gamma_{z}\gamma_{5}$ & Longitudinal-axial diquark \\
\bottomrule
\end{tabular}
\end{center}

In PyQUDA code:
\begin{verbatim}
Cg5  = (1j * gamma.gamma(2) @ gamma.gamma(8)) @ gamma.gamma(15)
CgT5 = (1j * gamma.gamma(2) @ gamma.gamma(8)) @ gamma.gamma(7)
CgZ5 = (1j * gamma.gamma(2) @ gamma.gamma(8)) @ gamma.gamma(11)
\end{verbatim}
Note that $\Cmat = i\gamma_{2}\gamma_{4}$ is built first,
then right-multiplied by the appropriate gamma structure.

\subsection{Epsilon Tensor Convention}

The color epsilon tensor is defined as
\begin{equation}
  \eps_{abc} = \begin{cases}
    +1 & \text{if }(a,b,c)\text{ is an even permutation of }(0,1,2),\\
    -1 & \text{if }(a,b,c)\text{ is an odd permutation of }(0,1,2),\\
    0  & \text{otherwise}.
  \end{cases}
\end{equation}
In the code, $\eps$ is initialized on the host as
\begin{verbatim}
epsilon_host = xp.zeros((3,3,3), dtype=...)
for a in range(3):
    b = (a+1) % 3
    c = (a+2) % 3
    epsilon_host[a,b,c] = 1
    epsilon_host[a,c,b] = -1
\end{verbatim}
This implements $\eps_{012} = \eps_{120} = \eps_{201} = +1$ and
$\eps_{021} = \eps_{102} = \eps_{210} = -1$, i.e., the standard convention
with cyclic permutations positive.

% ---------------------------------------------------------------------------
\section{Spin and Parity Projection}
% ---------------------------------------------------------------------------

The proton two- and three-point functions are projected onto definite parity and
spin states by a set of spin projection matrices $\PolProj$.  In the rest frame
of the proton, the relevant matrices constructed in the code are:

\subsection{Positive Parity Projector}

The positive parity projector isolates the forward-propagating (positive-energy)
component of the proton two-point function:
\begin{equation}
  \PolProj_{P} = \frac{\gamma_{0} + \gamma_{4}}{4}
               = \frac{1}{4}\big(\texttt{gamma(0)} + \texttt{gamma(8)}\big).
\end{equation}
Here $\gamma_{0} \equiv \mathbb{1}_{4\times 4}$ and $\gamma_{4}$ is the temporal
Dirac matrix.  In Euclidean space, this projector selects states with
$p_{4} \approx +im_{N}$.

\subsection{Spin Projectors}

For a spin-$\tfrac{1}{2}$ particle at rest, the spin projection operators in the
$z$ and $x$ directions are
\begin{align}
  \PolProj_{Sz}^{+} &= \gamma_{0} - i\gamma_{1}\gamma_{2}
                     = \texttt{gamma(0)} - i\,\texttt{gamma(1)}\,\texttt{@}\,\texttt{gamma(2)},
  \label{eq:Szp} \\
  \PolProj_{Sz}^{-} &= \gamma_{0} + i\gamma_{1}\gamma_{2}
                     = \texttt{gamma(0)} + i\,\texttt{gamma(1)}\,\texttt{@}\,\texttt{gamma(2)},
  \label{eq:Szm} \\
  \PolProj_{Sx}^{+} &= \gamma_{0} - i\gamma_{2}\gamma_{4}
                     = \texttt{gamma(0)} - i\,\texttt{gamma(2)}\,\texttt{@}\,\texttt{gamma(4)},
  \label{eq:Sxp} \\
  \PolProj_{Sx}^{-} &= \gamma_{0} + i\gamma_{2}\gamma_{4}
                     = \texttt{gamma(0)} + i\,\texttt{gamma(2)}\,\texttt{@}\,\texttt{gamma(4)}.
  \label{eq:Sxm}
\end{align}

\subsection{Combined Parity-Spin Projectors}
\label{sec:spinproj}

The physical proton states are selected by combining the parity and spin
projectors.  The \texttt{PolProjections} dictionary contains
\begin{align}
  \texttt{"PpSzp"}    &= \PolProj_{P}\,\PolProj_{Sz}^{+}, \\
  \texttt{"PpSzm"}    &= \PolProj_{P}\,\PolProj_{Sz}^{-}, \\
  \texttt{"PpSxp"}    &= \PolProj_{P}\,\PolProj_{Sx}^{+}, \\
  \texttt{"PpSxm"}    &= \PolProj_{P}\,\PolProj_{Sx}^{-}, \\
  \texttt{"PpUnpol"}  &= \PolProj_{P} \qquad\text{(unpolarized)}.
\end{align}
These five projectors are labeled by the \texttt{pol\_list} parameter and applied
as spin matrices with indices $(i,j) = (\text{spin\_sink}, \text{spin\_src})$.

% ---------------------------------------------------------------------------
\section{Proton Two-Point Function}
% ---------------------------------------------------------------------------

\subsection{Definition}

The connected proton two-point function with sink bilinear $\Gsink$ and momentum
projection $p$ is
\begin{equation}
  C_{2}^{\Gsink}(p,\,t_{\text{sep}})
  = \sum_{\mathbf{x}} \exp\!\big(-i\mathbf{p}\cdot(\mathbf{x}-\mathbf{x}_{0})\big)\;
    \PolProj_{\alpha\alpha'}\;
    \big\langle
      \chi_{\alpha}(\mathbf{x},\,t_{\text{sep}}+t_{0})\;
      {\Gsink}_{\beta\beta'}\;
      \bar{\chi}_{\beta'}(\mathbf{x}_{0},\,t_{0})
    \big\rangle,
  \label{eq:C2_def}
\end{equation}
where the source is placed at $\xsrc = (\mathbf{x}_{0}, t_{0})$ and the sink is at
$(\mathbf{x},\,t = t_{0}+t_{\text{sep}})$.  The sink gamma matrix $\Gsink$ is
scanned over all 16 bilinear structures listed in Section~\ref{sec:gammas}; in the
code they are indexed by the integer label $g \in \{0,\dots,15\}$.

The $\bar{\chi}$ conjugate field is obtained via
$\bar{\chi} = \chi^{\dagger}\gamma_{4}$.  After inserting the three-quark
interpolators~\eqref{eq:chi}, the two-point function contains six fermion fields
(three quarks and three antiquarks) whose Wick contractions produce the baryon
propagator.

\subsection{Wick Contraction Result}

Performing all Wick contractions of~\eqref{eq:C2_def} yields two topologically
distinct terms, corresponding to the two ways of pairing the identical up-quark
fields.  In the isospin-symmetric limit (identical propagators for $u$ and $d$
quarks), we denote the common quark propagator from $\xsrc$ to $x$ as
\begin{equation}
  \fwprop^{i,k}_{a,d}(x,\xsrc)
  \equiv \langle u^{i}_{a}(x)\,\bar{u}^{k}_{d}(\xsrc)\rangle
      = \langle d^{i}_{a}(x)\,\bar{d}^{k}_{d}(\xsrc)\rangle,
\end{equation}
where $(i,k)$ are spin\_sink and spin\_src indices and $(a,d)$ are color\_sink and
color\_src indices.

The two terms read:

\medskip\noindent
\textbf{Term 1 (direct contraction):}
\begin{equation}
  \boxed{
  \begin{aligned}
    C_{2}^{(1)}[g,p,t] = \sum_{\mathbf{x}} \Phi_{p}(\mathbf{x}-\mathbf{x}_{0})\;
    \eps_{abc}\,\eps_{def}\;
    (\Cmat\Gint)_{ij}\;(\Cmat\Gint)_{kl}\;
    [{\Gsink}_{g}]_{mn}\;
    &\fwprop^{i,k}_{a,d}(x,\xsrc)\;
    \fwprop^{j,l}_{b,e}(x,\xsrc)\;
    \fwprop^{m,n}_{c,f}(x,\xsrc).
  \end{aligned}
  }
  \label{eq:term1}
\end{equation}

\medskip\noindent
\textbf{Term 2 (exchange contraction):}
\begin{equation}
  \boxed{
  \begin{aligned}
    C_{2}^{(2)}[g,p,t] = \sum_{\mathbf{x}} \Phi_{p}(\mathbf{x}-\mathbf{x}_{0})\;
    \eps_{abc}\,\eps_{def}\;
    (\Cmat\Gint)_{ij}\;(\Cmat\Gint)_{kl}\;
    [{\Gsink}_{g}]_{mn}\;
    &\fwprop^{i,k}_{a,d}(x,\xsrc)\;
    \fwprop^{j,n}_{b,e}(x,\xsrc)\;
    \fwprop^{m,l}_{c,f}(x,\xsrc).
  \end{aligned}
  }
  \label{eq:term2}
\end{equation}

\subsection{Sign Convention}

The total two-point function is
\begin{equation}
  \boxed{
  C_{2}^{\Gsink}(p,t) = -\big(C_{2}^{(1)} + C_{2}^{(2)}\big).
  }
  \label{eq:C2_sign}
\end{equation}
The overall minus sign arises from the Grassmann nature of the fermion fields and
the conventions of the epsilon tensors.  It is essential for obtaining the
correct positive-definite spectral decomposition of the proton two-point function.

\subsection{Physical Interpretation of the Two Terms}

In Term~1 (\ref{eq:term1}), the index structure
$\fwprop^{i,k}_{a,d}\;\fwprop^{j,l}_{b,e}\;\fwprop^{m,n}_{c,f}$ corresponds to the
``direct'' pairing where the first up quark at the source connects to the first up
quark at the sink, the down quark connects to the down quark, and the second up
quark connects to the second up quark.  The spin indices contract as:
\begin{itemize}
  \item $(\Cmat\Gint)_{ij}$ couples the spin of propagators 1 and 2 (the diquark);
  \item $(\Cmat\Gint)_{kl}$ couples the spin of the conjugate diquark;
  \item $[{\Gsink}_{g}]_{mn}$ provides the sink spin structure.
\end{itemize}

Term~2 (\ref{eq:term2}) is the exchange term where the spin\_src indices $n$ and
$l$ are swapped between propagators 2 and 3.  This term is mandatory for fermionic
statistics: the two up quarks are identical, so their propagators must be
antisymmetrized.  The sum $C_{2}^{(1)}+C_{2}^{(2)}$ is antisymmetric under exchange
of the two up-quark propagators, as required by the Pauli principle.

% ---------------------------------------------------------------------------
\section{Optimized Two-Point Contraction}
% ---------------------------------------------------------------------------

The direct evaluation of~\eqref{eq:term1} and~\eqref{eq:term2} via a single einsum
would require a 10-index tensor contraction with memory footprint
$\mathcal{O}(V \times 4^{6} \times 3^{6})$, which is impractical.  The code in
\texttt{contract\_2pt\_TMD} splits the large contraction into a sequence of
smaller two- and three-tensor operations.  We detail each step below.

\subsection{Term 1: Step-by-Step Decomposition}

The einsum indices used in the code are:
\begin{itemize}
  \item \texttt{wtzyx} -- local spacetime coordinates (complex, lattice-wide);
  \item $i,j,k,l,m,n$ -- spin indices ($0,\dots,3$);
  \item $a,b,c,d,e,f$ -- color indices ($0,1,2$);
  \item $g$ -- sink gamma index ($0,\dots,15$);
  \item $p$ -- momentum index.
\end{itemize}

The full Term~1 contraction before optimization is
\begin{equation}
  \texttt{einsum("abc, def, ij, kl, gmn, p..., ...ikad, ...jlbe, ...mncf -> gpt")}.
\end{equation}

The code decomposes this into eight steps:

\medskip\noindent
\textbf{Step 1 (t1\_s1):} Contract $\eps_{abc}$ with the first propagator on
color index $a$:
\begin{equation}
  \texttt{t1\_s1 = einsum("abc, wtzyxikad -> wtzyxikbcd")},
\end{equation}
producing an intermediate of shape $[\texttt{wtzyx},i,k,b,c,d]$.  The indices
$(b,c)$ remain free for later epsilon contraction; $(d)$ is the color\_src index
that will contract with $\eps_{def}$.

\medskip\noindent
\textbf{Step 2 (t1\_s2):} Contract $(\Cmat\Gint)_{ij}$ with the second propagator
on spin index $j$:
\begin{equation}
  \texttt{t1\_s2 = einsum("ij, wtzyxjlbe -> wtzyxilbe")},
\end{equation}
producing $[\texttt{wtzyx},i,l,b,e]$.  Spin index $j$ is summed; $i$ will later
be contracted with t1\_s1.

\medskip\noindent
\textbf{Step 3 (term1\_sink):} Combine the two partial sink blocks by contracting
spin index $i$ and color index $b$:
\begin{equation}
  \texttt{term1\_sink = einsum("wtzyxikbcd, wtzyxilbe -> wtzyxklcde")},
\end{equation}
producing $[\texttt{wtzyx},k,l,c,d,e]$.  This object encodes the partially
assembled sink-side baryon contraction.

\medskip\noindent
\textbf{Step 4 (term1\_p3):} Contract the sink gamma matrix with the third
propagator on spin indices $(m,n)$:
\begin{equation}
  \texttt{term1\_p3 = einsum("gmn, wtzyxmncf -> gwtzyxcf")},
\end{equation}
producing $[g,\texttt{wtzyx},c,f]$.  This step also introduces the gamma index
$g$ for scanning all 16 sink bilinears.

\medskip\noindent
\textbf{Step 5 (t1\_f1):} Contract $\eps_{def}$ with term1\_sink on color indices
$(d,e)$:
\begin{equation}
  \texttt{t1\_f1 = einsum("def, wtzyxklcde -> wtzyxklcf")},
\end{equation}
producing $[\texttt{wtzyx},k,l,c,f]$.

\medskip\noindent
\textbf{Step 6 (t1\_f2):} Contract $(\Cmat\Gint)_{kl}$ to close the spin structure
of the diquark:
\begin{equation}
  \texttt{t1\_f2 = einsum("kl, wtzyxklcf -> wtzyxcf")},
\end{equation}
producing $[\texttt{wtzyx},c,f]$.

\medskip\noindent
\textbf{Step 7 (t1\_f3):} Contract with the P3 block on the remaining color
indices $(c,f)$:
\begin{equation}
  \texttt{t1\_f3 = einsum("wtzyxcf, gwtzyxcf -> gwtzyx")},
\end{equation}
producing $[g,\texttt{wtzyx}]$.

\medskip\noindent
\textbf{Step 8 (term1):} Contract with the momentum phases to project to definite
spatial momentum and sum over space:
\begin{equation}
  \texttt{term1 = einsum("pwtzyx, gwtzyx -> gpt")},
\end{equation}
producing the final shape $[g,p,t]$.

\subsection{Term 2: Step-by-Step Decomposition}

The full Term~2 contraction before optimization is
\begin{equation}
  \texttt{einsum("abc, def, ij, kl, gmn, p..., ...ikad, ...jnbe, ...mlcf -> gpt")}.
\end{equation}
Note the index differences from Term~1:
\begin{itemize}
  \item Prop 2 has indices $(\dots,j,n,b,e)$ instead of $(\dots,j,l,b,e)$;
  \item Prop 3 has indices $(\dots,m,l,c,f)$ instead of $(\dots,m,n,c,f)$.
\end{itemize}

The eight optimized steps for Term~2 are:

\medskip\noindent
\textbf{Step 1 (t2\_s1):} Same as t1\_s1 (eps on first propagator):
\begin{equation}
  \texttt{t2\_s1 = einsum("abc, wtzyxikad -> wtzyxikbcd")}.
\end{equation}

\medskip\noindent
\textbf{Step 2 (t2\_s2):} Contract $(\Cmat\Gint)_{ij}$ with the second propagator
on spin index $j$ (note: $n$ remains free instead of $l$):
\begin{equation}
  \texttt{t2\_s2 = einsum("ij, wtzyxjnbe -> wtzyxinbe")},
\end{equation}
producing $[\texttt{wtzyx},i,n,b,e]$.

\medskip\noindent
\textbf{Step 3 (term2\_sink):} Combine on indices $(i,b)$, keeping $k,n$ free:
\begin{equation}
  \texttt{term2\_sink = einsum("wtzyxikbcd, wtzyxinbe -> wtzyxkncde")},
\end{equation}
producing $[\texttt{wtzyx},k,n,c,d,e]$.

\medskip\noindent
\textbf{Step 4 (term2\_p3):} Contract sink gamma with third propagator (note:
$n,l$ both free):
\begin{equation}
  \texttt{term2\_p3 = einsum("gmn, wtzyxmlcf -> gwtzyxnlcf")},
\end{equation}
producing $[g,\texttt{wtzyx},n,l,c,f]$.

\medskip\noindent
\textbf{Step 5 (t2\_f1):} Contract $\eps_{def}$ with term2\_sink:
\begin{equation}
  \texttt{t2\_f1 = einsum("def, wtzyxkncde -> wtzyxkncf")},
\end{equation}
producing $[\texttt{wtzyx},k,n,c,f]$.

\medskip\noindent
\textbf{Step 6 (t2\_f2):} Contract t2\_f1 with term2\_p3 on indices $(n,c,f)$:
\begin{equation}
  \texttt{t2\_f2 = einsum("wtzyxkncf, gwtzyxnlcf -> gwtzyxkl")},
\end{equation}
producing $[g,\texttt{wtzyx},k,l]$.

\medskip\noindent
\textbf{Step 7 (t2\_f3):} Contract $(\Cmat\Gint)_{kl}$:
\begin{equation}
  \texttt{t2\_f3 = einsum("kl, gwtzyxkl -> gwtzyx")},
\end{equation}
producing $[g,\texttt{wtzyx}]$.

\medskip\noindent
\textbf{Step 8 (term2):} Momentum projection:
\begin{equation}
  \texttt{term2 = einsum("pwtzyx, gwtzyx -> gpt")}.
\end{equation}

\subsection{Memory Optimization Rationale}

The key insight of the split contraction is that the full six-propagator tensor
(three propagators $\times$ two spin indices $\times$ two color indices each) is
never materialized.  Instead:
\begin{itemize}
  \item Steps 1--3 build the sink block using only \emph{two} propagators at a time,
        with intermediate tensors of size $\mathcal{O}(V \times 4^{3} \times 3^{2})$.
  \item Step 4 processes the third propagator with the sink gamma independently.
  \item Steps 5--8 assemble the final result using tensors where most spin indices
        have already been contracted.
\end{itemize}
The largest intermediate tensor is \texttt{term1\_sink} or \texttt{term2\_sink}
with shape $[\texttt{wtzyx}, 4, 4, 3, 3, 3]$, which is
$V \times 4^{2} \times 3^{3} = 432\,V$ complex numbers -- manageable for typical
lattice volumes.

% ===========================================================================
\part{Proton Sequential Source and Connected Three-Point Function}
% ===========================================================================

% ---------------------------------------------------------------------------
\section{Connected Proton Three-Point Function}
% ---------------------------------------------------------------------------

\subsection{Definition before the Sequential Trick}
\label{sec:C3_def}

The connected proton three-point function with a local insertion $O_{f}(x)$ on
flavor $f \in \{U,D\}$ is
\begin{equation}
  \boxed{
  \begin{aligned}
  C_{3}^{f}(\momtrans,\tau;\,\pfinal,t_{\text{sep}},\PolProj)
  = \sum_{\mathbf{x},\mathbf{y}}\;
    &\Phi_{\momtrans}(\xins-\xsrc)\;
    \Phi_{\pfinal}(\xsink-\xsrc)\;
    \PolProj_{\alpha\alpha'}\\
    &\times\big\langle
        \chi_{\alpha}(\xsink)\;
        O_{f}(\xins)\;
        \bar{\chi}_{\alpha'}(\xsrc)
      \big\rangle_{\text{connected}}.
  \end{aligned}
  }
  \label{eq:C3_def}
\end{equation}
Here:
\begin{itemize}
  \item $\xsrc = (\mathbf{x}_{0},t_{0})$ is the source position;
  \item $\xins = (\mathbf{x},\tau)$ is the EMT insertion point;
  \item $\xsink = (\mathbf{y},\,t_{0}+t_{\text{sep}})$ is the fixed proton sink;
  \item $\momtrans = \pfinal - \pini$ is the momentum transfer;
  \item $\pfinal$ is the final-state proton momentum (set by \texttt{pf});
  \item $t_{\text{sep}}$ is the source-sink separation;
  \item $\PolProj$ is one of the five combined spin-parity projectors from
        Section~\ref{sec:spinproj}.
\end{itemize}
The subscript ``connected'' means that all three valence quark lines connect the
source, insertion, and sink -- no disconnected (sea-quark-loop) diagrams are
included at this stage.

\subsection{Sequential-Source Strategy}

The double sum over sink position $\mathbf{y}$ and insertion position $\mathbf{x}$
cannot be evaluated by direct lattice summation for each operator insertion, since
it would require a new inversion for every $(t_{\text{sep}}, \mathbf{x})$.
Instead, the fixed-sink sequential-source method is employed:

\begin{enumerate}
  \item Build the forward (source-to-sink) propagator $\fwprop(\xsink,\xsrc)$
        and apply sink smearing.
  \item Contract the sink-side baryon (diquark + spectator) to produce a
        \emph{sequential source} $\seqsrc(x;\,\pfinal,t_{\text{sep}},\PolProj)$.
  \item Invert the Dirac operator on the sequential source:
        $D\,\seqprop = \seqsrc$.
  \item The sequential propagator $\seqprop$ carries the full sink-side
        information.  The three-point function is then obtained by a single
        contraction of $\seqprop$ with the forward propagator and the operator
        insertion at each $x$.
\end{enumerate}

This reduces the computational cost from $\mathcal{O}(N_{t_{\text{sep}}} \times V)$
inversions to $\mathcal{O}(N_{t_{\text{sep}}} \times N_{\text{flavor}} \times
N_{\text{pol}})$ inversions, where $N_{\text{pol}}$ is the number of polarization
states.

% ---------------------------------------------------------------------------
\section{Proton Fixed-Sink Sequential Source}
% ---------------------------------------------------------------------------

The sequential source is built by \texttt{create\_bw\_seq\_pyquda}, which
delegates the flavor-specific diquark contraction to either
\texttt{up\_quark\_insertion\_pyquda} (flavor $U$, insertion on an up quark) or
\texttt{down\_quark\_insertion\_pyquda} (flavor $D$, insertion on the down quark).

\subsection{Flavor 1: Up-Quark Insertion}

When the operator insertion is on an up-quark line, the proton contains two
identical up quarks, which leads to four distinct contraction terms after
antisymmetrization.  We define the following intermediate objects built from
the forward propagator $\fwprop$ (using isospin symmetry, $\fwprop_{u}=\fwprop_{d}=\fwprop$):

\begin{align}
  [\Gamma^{T}\!\fwprop\,\Gamma]^{ij}_{ab}
    &\equiv (\Gint^{T})_{ik}\;\fwprop^{kl}_{ab}\;(\Gint)_{lj},
    \label{eq:GtDG} \\
  [\PolProj\,\fwprop]^{ij}_{ab}
    &\equiv \PolProj_{ik}\;\fwprop^{kj}_{ab},
    \label{eq:PDu} \\
  [\fwprop\,\PolProj]^{ij}_{ab}
    &\equiv \fwprop^{ik}_{ab}\;\PolProj_{kj},
    \label{eq:DuP} \\
  [\Tr_{\text{spin}}(\fwprop\,\PolProj)]_{ab}
    &\equiv \fwprop^{kj}_{ab}\;\PolProj_{jk}.
    \label{eq:TrDuP}
\end{align}
Equation~\eqref{eq:GtDG} is the down-quark propagator dressed with the
interpolating gamma on both sides -- this is the standard structure appearing
in the diquark.  Equations~\eqref{eq:PDu} and~\eqref{eq:DuP} are the up-quark
propagator multiplied by the spin projector from the left or right.
Equation~\eqref{eq:TrDuP} is the spin-space trace of $\fwprop\PolProj$, leaving
only color indices.

\medskip\noindent
With these definitions, the \textbf{four terms} of the up-quark sequential source
are:

\medskip\noindent
\textbf{Term R1:} Both up quarks participate in the diquark; the spin projector
is applied to the output:
\begin{equation}
  \boxed{
  R_{1}^{ij;\,cf} = \PolProj^{ij}\;
  \eps_{abc}\,\eps_{def}\;
  \Big[(\Gamma^{T}\fwprop\Gamma)^{mn}_{be}\;
        \fwprop^{mn}_{ad}\Big],
  }
  \label{eq:R1}
\end{equation}
where the spin indices $m,n$ are summed (the term $[\Gamma^{T}\fwprop\Gamma \cdot
\fwprop]$ is a Lorentz scalar in spin space).  The result carries two free spin
indices $(i,j)$ from the projector and two free color indices $(c,f)$ from the
epsilon tensors.

\medskip\noindent
\textbf{Term R2:} The spin trace of $\fwprop\PolProj$ is used, and the diquark
$\Gamma^{T}\fwprop\Gamma$ provides the spin structure:
\begin{equation}
  \boxed{
  R_{2}^{ij;\,cf} =
  \eps_{abc}\,\eps_{def}\;
  [\Tr_{\text{spin}}(\fwprop\,\PolProj)]_{ad}\;
  [\Gamma^{T}\fwprop\Gamma]^{ji}_{be},
  }
  \label{eq:R2}
\end{equation}
where the spin indices $(j,i)$ of $\Gamma^{T}\fwprop\Gamma$ become the output
spin indices $(i,j)$.

\medskip\noindent
\textbf{Term R3:} Spin contraction between $[\PolProj\,\fwprop]$ and
$[\Gamma^{T}\fwprop\Gamma]$:
\begin{equation}
  \boxed{
  R_{3}^{ij;\,cf} =
  \eps_{abc}\,\eps_{def}\;
  [\PolProj\,\fwprop]^{ik}_{ad}\;
  [\Gamma^{T}\fwprop\Gamma]^{jk}_{be},
  }
  \label{eq:R3}
\end{equation}
where the repeated spin index $k$ is summed.

\medskip\noindent
\textbf{Term R4:} Spin contraction between $[\Gamma^{T}\fwprop\Gamma]$ and
$[\fwprop\,\PolProj]$:
\begin{equation}
  \boxed{
  R_{4}^{ij;\,cf} =
  \eps_{abc}\,\eps_{def}\;
  [\Gamma^{T}\fwprop\Gamma]^{ki}_{ad}\;
  [\fwprop\,\PolProj]^{kj}_{be},
  }
  \label{eq:R4}
\end{equation}
where again the spin index $k$ is summed.

\medskip\noindent
The total up-quark sequential source is
\begin{equation}
  \boxed{
  D_{\text{total}}^{ij;\,cf} = -\big(R_{1} + R_{2} + R_{3} + R_{4}\big).
  }
  \label{eq:Dtotal}
\end{equation}
The overall minus sign is consistent with the minus sign in the two-point
function~\eqref{eq:C2_sign}, as verified by the fact that the sequential
propagator contracted with the forward propagator reproduces the two-point
function when the insertion is the identity operator.

\subsection{Flavor 2: Down-Quark Insertion}

For the down-quark insertion there is only one down quark in the proton, so the
antisymmetrization produces only two terms.  The intermediate objects are
\begin{align}
  [\PolProj\,\fwprop]^{fc} &\equiv \PolProj_{fk}\,\fwprop^{k}_{c}
    \quad\text{(spin\_sink $f$, spin\_src $c$)}, \label{eq:PDu_down} \\
  [\Gamma^{T}\fwprop\Gamma]^{uv}_{eb}
    &\equiv (\Gint^{T})_{uk}\,\fwprop^{kl}_{eb}\,{\Gint}_{lv},
    \label{eq:GtDG_down} \\
  [\Gamma^{T}\fwprop]^{uj}_{ec}
    &\equiv (\Gint^{T})_{uk}\,\fwprop^{kj}_{ec},
    \label{eq:GtD_down} \\
  [\PolProj\,\fwprop\,\Gamma]^{jk}_{fb}
    &\equiv \PolProj_{jl}\,\fwprop^{lm}_{fb}\,{\Gint}_{mk}.
    \label{eq:PDG_down}
\end{align}
(For compactness we suppress the common spacetime and color indices in the
notation above; the full forms are given in the einsum expressions below.)

\medskip\noindent
The two terms are:

\medskip\noindent
\textbf{Term 1 (down insertion):}
\begin{equation}
  \boxed{
  {D_{1}}^{uv;\,ad}_{\text{(down)}} =
  \eps_{abc}\,\eps_{def}\;
  [\PolProj\,\fwprop]^{fc}\;
  [\Gamma^{T}\fwprop\Gamma]^{uv}_{eb},
  }
  \label{eq:Ddown1}
\end{equation}
where the spin structure comes entirely from $\Gamma^{T}\fwprop\Gamma$, and
$[\PolProj\,\fwprop]$ is traced in all internal indices that connect to the
epsilon contraction (the spin index $c$ and color indices are summed through the
epsilons).

\medskip\noindent
\textbf{Term 2 (down insertion):}
\begin{equation}
  \boxed{
  {D_{2}}^{uv;\,ad}_{\text{(down)}} =
  \eps_{abc}\,\eps_{def}\;
  [\Gamma^{T}\fwprop]^{uj}_{ec}\;
  [\PolProj\,\fwprop\,\Gamma]^{jk}_{fb},
  }
  \label{eq:Ddown2}
\end{equation}
where the spin index $j$ and color index $f$ are summed with the epsilon
contraction.

\medskip\noindent
The total down-quark sequential source is
\begin{equation}
  \boxed{
  D_{\text{(down)}}^{uv;\,ad} = D_{2} - D_{1}.
  }
  \label{eq:Ddown_total}
\end{equation}

\subsection{From Diquark Contraction to Sequential Inversion}

After the diquark contraction, the sequential source construction proceeds as
follows (for each flavor and polarization):

\begin{enumerate}
  \item \textbf{Time slicing.}  The diquark-contracted object is restricted to
        the sink time slice $t_{\text{sink}} = (t_{0} + t_{\text{sep}}) \bmod L_{t}$
        using \texttt{sequential12}, which sets the source to zero at all other
        time slices.

  \item \textbf{Momentum phase and $\gamma_{5}$ conjugation.}  The sliced source
        is multiplied by the final-state momentum phase
        $\Phi_{\pfinal}(\xsink-\xsrc)$ and conjugated with $\gamma_{5}$:
        \begin{equation}
          {\seqsrc}^{i,k}_{a,b}(\xsink)
          = [\gamma_{5}]_{ij}\;
            \Phi_{\pfinal}(\xsink-\xsrc)\;
            [D^{\dagger}]^{j,k}_{b,a},
        \end{equation}
        where $D$ is the diquark contraction result from~\eqref{eq:Dtotal} or
        \eqref{eq:Ddown_total}.  This $\gamma_{5}$-hermiticity step is required
        to convert the source into the correct form for the sequential inversion
        with the same Dirac operator used for the forward propagator.

  \item \textbf{Boosted smearing.}  The sequential source is smeared using
        \texttt{boosted\_smearing} with width $w$ and boost vector
        \texttt{boost\_out}.  This applies the same gauge-covariant Gaussian
        smearing used at the source and sink of the two-point function.

  \item \textbf{Sequential inversion.}  The smeared source is inverted:
        \begin{equation}
          \seqprop = D^{-1}\,{\seqsrc}_{\text{smeared}}.
        \end{equation}

  \item \textbf{Final $\gamma_{5}$ conjugation.}  The sequential propagator is
        conjugated once more with $\gamma_{5}$ to obtain the final object:
        \begin{equation}
          {\seqprop_{\text{final}}}^{i,j}_{c,f}(x)
          = [\seqprop^{\dagger}]^{j,i}_{f,c}(x)\;[\gamma_{5}]_{i,j}.
        \end{equation}
        In einsum notation:
        \begin{equation}
          \texttt{final\_term = einsum("wtzyxijfc, ik -> wtzyxjkcf",}
          \;\texttt{seq\_prop.conj(), G5)}.
        \end{equation}
\end{enumerate}

\subsection{Sequential Object Shape}

The final sequential object (one per polarization) has the shape
\begin{equation}
  \boxed{
  \seqprop[\PolProj]\; \longleftrightarrow\;
  \texttt{[w,\, t,\, z,\, y,\, x\_cb,\, spin\_a,\, spin\_b,\, color\_a,\, color\_b]},
  }
\end{equation}
where:
\begin{itemize}
  \item \texttt{spin\_a}, \texttt{spin\_b} are the Dirac indices at the two ends
        of the sequential propagator (corresponding to the insertion and sink ends);
  \item \texttt{color\_a}, \texttt{color\_b} are the corresponding color indices.
\end{itemize}
In einsum notation we label these as indices $(j,i,c,f)$ for
(\texttt{spin\_a}, \texttt{spin\_b}, \texttt{color\_a}, \texttt{color\_b}).
Note the index ordering: the sequential propagator's ``source'' (insertion-side)
indices are $(i,f)$ and its ``sink''-side indices are $(j,c)$.

The full list of sequential objects is gathered into an array
\begin{equation}
  \texttt{seq\_bw}\;\longleftrightarrow\;
  \texttt{[polarization,\, w,\, t,\, z,\, y,\, x\_cb,\, spin\_a,\, spin\_b,\, color\_a,\, color\_b]}.
\end{equation}

% ===========================================================================
\part{Proton EMT Contractions}
% ===========================================================================

% ---------------------------------------------------------------------------
\section{Scalar Diagnostic: $C_{3}^{\chi}$}
% ---------------------------------------------------------------------------

The scalar diagnostic three-point function tests the sequential-source machinery
by contracting the sequential propagator with the forward propagator without any
derivative insertion.  For flavor $f$,
\begin{equation}
  \boxed{
  C_{3}^{\chi,\,f}(\momtrans,\tau)
  = \sum_{\mathbf{x}}
    \Phi_{\momtrans}(\xins-\xsrc)\;
    \big[\seqprop^{f}(x)\big]^{j,i}_{c,f}\;
    \big[\fwprop(x,\xsrc)\big]^{i,j}_{f,c},
  }
  \label{eq:C3_chi}
\end{equation}
where the spin indices $(j,i)$ of the sequential propagator are contracted with
$(i,j)$ of the forward propagator, and color indices $(c,f)$ with $(f,c)$.  This
is the full spin-color trace that, for $\momtrans=0$, reproduces the zero-momentum
two-point function $C_{2}(t_{\text{sep}})$ up to a normalization factor.

In einsum notation:
\begin{equation}
  \texttt{C3\_chi = contract("wtzyxjicf, wtzyxijfc -> wtzyx", seq\_data, prop\_fw.data)}.
\end{equation}
The resulting scalar field is then projected to spatial momentum and the time
dependence is extracted.

% ---------------------------------------------------------------------------
\section{Connected Quark EMT Insertion}
% ---------------------------------------------------------------------------

\subsection{The Quark EMT Operator}

In Euclidean continuum QCD, the symmetric traceless quark energy-momentum tensor is
\begin{equation}
  \TmunuQ(x) = \frac{1}{4}\,
    \bar{\psi}(x)\Big(
      \gamma_{\mu}\overleftrightarrow{D}_{\nu}
      + \gamma_{\nu}\overleftrightarrow{D}_{\mu}
    \Big)\psi(x),
\end{equation}
where $\overleftrightarrow{D}_{\mu} = \overrightarrow{D}_{\mu} -
\overleftarrow{D}_{\mu}$ is the symmetric covariant derivative.  On the lattice, we
implement this as a sum of two terms: the derivative acting on the forward quark
line and the derivative acting on the sequential (backward) quark line.

\subsection{First Term: Derivative on the Forward Line}

The forward quark line runs from the source to the insertion point.  Applying the
covariant derivative to it gives
\begin{equation}
  \boxed{
  C_{3,\,f,\,\mu\nu}^{\text{first}}(\momtrans,\tau)
  = +\frac{1}{2}\sum_{\mathbf{x}}
    \Phi_{\momtrans}(\xins-\xsrc)\;
    \big[\seqprop^{f}(x)\big]^{j,i}_{c,f}\;
    \big[\gamma_{\nu}\,\Dmu\,\fwprop(x,\xsrc)\big]^{i,j}_{f,c},
  }
  \label{eq:C3_first}
\end{equation}
where the positive sign and the factor $1/2$ come from the decomposition of the
symmetrized EMT into forward and backward derivatives.

The implementation first applies the symmetric covariant derivative to the forward
propagator:
\begin{equation}
  \Dmu\fwprop = \frac{1}{2}\big(D_{+\mu} - D_{-\mu}\big)\,\fwprop,
\end{equation}
where $D_{\pm\mu}$ are the gauge-covariant forward/backward finite differences
with the appropriate gauge link parallel transport.  Then $\gamma_{\nu}$ is
contracted on the spin\_sink index:
\begin{equation}
  \texttt{gamma\_D\_fw = contract("ab, wtzyxbdij -> wtzyxadij", gamma\_nu, D\_fw.data)}.
\end{equation}
Finally, the spin-color trace with the sequential propagator is formed:
\begin{equation}
  \texttt{scalar\_field = 0.5 * contract("wtzyxjicf, wtzyxijfc -> wtzyx",
    seq\_data, gamma\_D\_fw)}.
\end{equation}

\subsection{Second Term: Derivative on the Sequential Line}

The sequential propagator runs from the sink to the insertion point (backward in
time).  The left-acting derivative is applied to it:
\begin{equation}
  \boxed{
  C_{3,\,f,\,\mu\nu}^{\text{second}}(\momtrans,\tau)
  = -\frac{1}{2}\sum_{\mathbf{x}}
    \Phi_{\momtrans}(\xins-\xsrc)\;
    \big[\overleftarrow{\Dmu}\,\seqprop^{f}(x)\big]^{j,i}_{c,f}\;
    \big[\gamma_{\nu}\,\fwprop(x,\xsrc)\big]^{i,j}_{f,c},
  }
  \label{eq:C3_second}
\end{equation}
where $\overleftarrow{\Dmu}$ denotes the covariant derivative acting to the left
(on the sequential propagator's source index).  The negative sign arises because
$\overleftarrow{D}_{\mu} = -\overrightarrow{D}_{\mu}^{\dagger}$ in the bilinear
form.

In the implementation, the left derivative is applied to the raw sequential
propagator before the final proton-specific $\gamma_{5}$ conjugation:
\begin{enumerate}
  \item \texttt{create\_bw\_seq\_raw\_pyquda(...)} returns only the raw
        sequential propagator immediately after the sequential inversion.  This
        avoids building and storing the finalized qTMD-style \texttt{dst\_seq}
        object in the proton EMT workflow.
  \item The symmetric covariant derivative $\frac{1}{2}(D_{+\mu}-D_{-\mu})$ is
        applied to this raw sequential propagator.
  \item The derivative result is converted to the final proton sequential layout
        with the same $\gamma_{5}$ hermiticity/index transformation used in
        \texttt{create\_bw\_seq\_pyquda}:
        \begin{equation}
          \texttt{leftD\_seq = contract("wtzyxijfc, ik -> wtzyxjkcf",}
          \;\texttt{D\_raw\_seq.conj(), G5)}.
        \end{equation}
\end{enumerate}
This mirrors the meson EMT construction, where the left derivative is built by
applying the covariant derivative to the raw sequential propagator and only then
forming the backward line through $\gamma_{5}$ hermiticity.  The convention is
also the lattice realization of the standard quark EMT operator used in proton
GFF calculations, $\bar\psi\,\overleftrightarrow{D}_{\{\mu}\gamma_{\nu\}}\psi$.
This is the same operator-level convention used, for example, in the
MIT/Fermilab proton gravitational form-factor calculation by Hackett, Pefkou,
and Shanahan, Phys. Rev. Lett. 132, 251904 (2024),
doi:10.1103/PhysRevLett.132.251904.

\subsection{Full Connected Quark EMT}
\label{sec:C3_total}

The total connected quark EMT three-point function is the sum of the two terms:
\begin{equation}
  \boxed{
  C_{3,\,f,\,\mu\nu}^{q}(\momtrans,\tau)
  = C_{3,\,f,\,\mu\nu}^{\text{first}}(\momtrans,\tau)
  + C_{3,\,f,\,\mu\nu}^{\text{second}}(\momtrans,\tau).
  }
  \label{eq:C3_total}
\end{equation}
After computing the full tensor, explicit symmetrization is enforced:
\begin{equation}
  C_{3,\,f,\,\mu\nu}^{q} \to \frac{1}{2}\big(C_{3,\,f,\,\mu\nu}^{q} + C_{3,\,f,\,\nu\mu}^{q}\big),
  \qquad
  C_{3,\,f,\,\nu\mu}^{q} = C_{3,\,f,\,\mu\nu}^{q}.
\end{equation}
This guarantees $T_{\mu\nu} = T_{\nu\mu}$ at the level of the measured data.

% ---------------------------------------------------------------------------
\section{Key Differences from the Meson EMT}
% ---------------------------------------------------------------------------

The proton EMT contraction differs from the meson (pion) case in several important
ways:

\begin{enumerate}
  \item \textbf{No \texttt{\_make\_dst2} required.}
        In the meson code, the backward sequential line is formed by
        $\gamma_{5}\,\seqprop^{\dagger}\,\gamma_{5}$ (the \texttt{dst2} object)
        before each contraction.  In the proton code, the sequential object
        \emph{already} has the correct spin structure built in during the diquark
        contraction and the final $\gamma_{5}$ conjugation step.  Therefore
        $\seqprop$ can be used directly without the \texttt{dst2} transformation.

  \item \textbf{Direct einsum contraction.}
        The meson contraction is a three-way trace:
        \begin{equation}
          \texttt{contract("wtzyxabij, wtzyxbcji, ca -> wtzyx", dst2, gamma\_D\_fw, src\_gamma)}.
        \end{equation}
        The proton contraction is simpler:
        \begin{equation}
          \texttt{contract("wtzyxjicf, wtzyxijfc -> wtzyx", seq\_data, prop\_fw\_data)},
        \end{equation}
        summing over spin indices $(j,i)$ and color indices $(c,f)$.  There is no
        source gamma matrix in the proton contraction because the spin projection
        is already absorbed into the sequential source.

  \item \textbf{Covariant derivative on the raw sequential propagator.}
        In the meson case, the left derivative is applied to the raw sequential
        propagator and then converted to the backward line through
        \texttt{\_left\_covdev\_dst2\_from\_prop}.  The proton EMT code now follows
        the same logic: the derivative acts on the raw sequential propagator, and
        the result is then converted to the final proton sequential layout through
        $\gamma_{5}$ hermiticity.  The already-finalized \texttt{seq\_data} is not
        treated as an ordinary \texttt{LatticePropagator} for the left derivative.

  \item \textbf{Flavor dimension.}  The proton result carries a flavor index
        $[U, D]$ since the insertion can be on either quark line.  The meson
        result has no flavor index (only one valence flavor).

  \item \textbf{Polarization dimension.}  The proton result has a polarization
        index from \texttt{pol\_list} (five entries).  The pion has no spin
        degree of freedom.

  \item \textbf{Time separation loop.}  The proton code loops over
        \texttt{t\_separations}, building new sequential sources for each
        $t_{\text{sep}}$.  The meson code has the same structure, but the proton
        sequential source construction involves the epsilon-tensor diquark
        contraction instead of a simple meson two-point contraction.
\end{enumerate}

% ---------------------------------------------------------------------------
\section{Gradient Flow}
% ---------------------------------------------------------------------------

\subsection{Flow Schedule}

The gradient flow is applied identically to the meson and proton codes.  The flow
schedule is:

\begin{center}
\begin{tabular}{ccl}
\toprule
Step index & Flow action & Description \\
\midrule
$s = 0$ & None (measure unflowed) & Unflowed (bare) operators \\
$s = 0 \to 1$ & 10 sub-steps of $\varepsilon/10$ each & First flow interval subdivided \\
$s = 1 \to 2$ & 1 step of $\varepsilon$ & Regular flow step \\
$\cdots$ & $\cdots$ & $\cdots$ \\
$s = N-1 \to N$ & 1 step of $\varepsilon$ & Final flow step \\
\bottomrule
\end{tabular}
\end{center}

Thus output index \texttt{step} corresponds to flow time $t_{\text{flow}} =
\texttt{step} \times \varepsilon$, with $t_{\text{flow}} = 0$ at
\texttt{step = 0}.  The subdivision of the first interval ensures that the initial
condition of the flow equation is well resolved numerically.

\subsection{Flowing Fermion and Gauge Fields Together}

At each flow step, both the forward propagator and the sequential propagator are
flowed simultaneously on the same flowed gauge background.  This is accomplished
by packing both propagators into a single \texttt{MultiLatticeFermion} and calling
\texttt{gradientFlow} once:
\begin{equation}
  (\fwprop_{\text{flowed}},\;\seqprop_{\text{flowed}})
  = \texttt{gradientFlow}\big([\fwprop,\seqprop],\;
    \texttt{flow\_type},\;N_{\text{steps}},\;\Delta t\big).
\end{equation}
The gauge field used for the fermion flow is set to anti-periodic temporal boundary
conditions via \texttt{U\_f.setAntiPeriodicT()}.  A copy of the unflowed gauge
field is used (\texttt{U\_f = U.copy()}) so that the original gauge field is
preserved for subsequent measurements.

\subsection{Physical Significance}

Gradient flow smooths ultraviolet fluctuations at a physical radius approximately
$\sqrt{8t_{\text{flow}}}$.  The EMT observables measured at finite flow time are
flowed composite operators.  Their flow-time dependence is analyzed using the
small-flow-time expansion to extract renormalized matrix elements.  This analysis
is performed \emph{outside} the contraction kernel; the kernel only produces the
flowed data at each \texttt{step}.

The fermion-flow kernel is four dimensional, so the quoted radius includes the
Euclidean time direction and is a diffusion scale rather than a strict support
boundary.  This matters for stochastic disconnected loops, as summarized in
Sec.~\ref{sec:proton_disc_full_volume_source}.

% ---------------------------------------------------------------------------
\section{Code Structure Overview}
% ---------------------------------------------------------------------------

The proton EMT measurement code is organized in the class
\texttt{ProtonQuarkEMT}, which inherits from
\texttt{EMTDisconnectedQuark1pt}.  The inheritance provides:
\begin{itemize}
  \item The stochastic quark one-point loop measurement
        (\texttt{flowed\_fermionic\_1pt});
  \item The covariant derivative utilities
        (\texttt{\_covdev\_sym\_prop}, \texttt{\_make\_dst2}, etc.);
  \item The propagator flow utilities (\texttt{\_advance\_flowed\_props});
  \item Gamma matrix helpers (\texttt{\_dirac\_gammas\_for},
        \texttt{\_gamma5\_for}, \texttt{\_get\_interpolator\_gamma\_for}).
\end{itemize}

The key methods of \texttt{ProtonQuarkEMT} are:

\medskip\noindent
\texttt{connected\_3pt(gauge, invPara, src\_pos, t\_separations, spin, tag, ...)}:
\begin{itemize}
  \item Builds the Dirac solver and the forward point-source propagator
        (via \texttt{\_make\_source\_prop}).
  \item Contracts the proton two-point function for normalization
        (via \texttt{contract\_proton\_2pt}, which reuses
        \texttt{proton\_TMD.contract\_2pt\_TMD}).
  \item Loops over $t_{\text{sep}}$ and flavor, building sequential sources
        via \texttt{create\_bw\_seq\_pyquda}.
  \item Loops over polarization and flow steps, calling
        \texttt{get\_C3\_chi\_proton} and
        \texttt{get\_C3\_Tmunu\_symmetrized\_proton} at each step.
  \item Advances the flow between steps.
  \item Saves results to HDF5.
\end{itemize}

\medskip\noindent
\texttt{get\_C3\_chi\_proton(latt\_info, prop\_fw, seq\_data, phases\_3pt, t0)}:
\begin{itemize}
  \item Computes the scalar diagnostic $C_{3}^{\chi}$ by direct einsum contraction.
\end{itemize}

\medskip\noindent
\texttt{get\_C3\_Tmunu\_symmetrized\_proton(U\_f, prop\_fw, seq\_data, phases\_3pt, t0)}:
\begin{itemize}
  \item Computes the full $4\times 4$ quark EMT three-point function.
  \item Loops over $\mu,\nu$ to evaluate the first-term (derivative on forward line)
        and second-term (derivative on sequential line) contributions.
  \item Symmetrizes the result.
\end{itemize}

\medskip\noindent
\texttt{\_covdev\_sym\_seq(U\_f, seq\_data, mu)}:
\begin{itemize}
  \item Wraps \texttt{seq\_data} as a \texttt{LatticePropagator}.
  \item Applies the symmetric covariant derivative using
        \texttt{\_covdev\_sym\_prop}.
  \item Returns the derivative data.
\end{itemize}

\medskip\noindent
\texttt{\_seq\_to\_prop(latt\_info, seq\_data)}:
\begin{itemize}
  \item Creates a \texttt{LatticePropagator} from the sequential data array
        for use with propagator-level methods.
\end{itemize}

% ---------------------------------------------------------------------------
\section{Output Structure}
% ---------------------------------------------------------------------------

\subsection{HDF5 Output for Connected Proton Three-Point}

The proton connected EMT three-point files save only three-point observables and
their momentum-transfer metadata.  The matching two-point functions are written
separately under the \texttt{EMTproton2pt} output tree.

\begin{center}
\begin{tabular}{lll}
\toprule
Dataset & Shape & Description \\
\midrule
\texttt{C3\_chi} & \texttt{[Nflavor, Npol, Nts, Nsteps+1, Nq, t\_sep+2]}
                  & Scalar diagnostic three-point function \\
\texttt{C3\_Tmunu} & \texttt{[Nflavor, Npol, Nts, Nsteps+1, Nq, 4, 4, t\_sep+2]}
                    & Quark EMT three-point function \\
\texttt{momentum\_transfer\_list} & \texttt{[Nq, 4]}
                                   & List of momentum transfer vectors \\
\bottomrule
\end{tabular}
\end{center}

The dimensions are:
\begin{itemize}
  \item $N_{\text{flavor}} = 2$: index 0 = up-quark insertion, index 1 = down-quark insertion;
  \item $N_{\text{pol}}$ = number of polarization states from \texttt{pol\_list} (typically 5);
  \item $N_{\text{ts}}$ = number of source-sink separations;
  \item $N_{\text{steps}}+1$ = number of flow times (including unflowed);
  \item $N_{q}$ = number of momentum transfers;
  \item $N_{t}$ = temporal lattice extent.
\end{itemize}

The attributes stored in the HDF5 file include the measurement type, spin,
flow parameters, source time, interpolator, flavor/polarization axis labels,
and the selected two-point momentum used for normalization.  Future files also
record non-destructive convention metadata:
\begin{itemize}
  \item \code{operator\_normalization = unringed\_flowed\_bilinear};
  \item \code{ringed\_normalization\_applied = False};
  \item \code{flow\_times = flow\_epsilon * [0,\ldots,Nsteps]};
  \item \code{quark\_flow\_scope = inserted\_operator\_quark\_legs\_only};
  \item \code{nucleon\_interpolator\_flowed = False};
  \item \code{ringed\_factor\_source = analysis\_from\_quark\_1pt\_kinetic}.
\end{itemize}
These attributes do not change any dataset or contraction.  They only make
explicit that \code{C3\_Tmunu} is an unringed flowed bilinear and must not be
treated as already ringed in downstream analysis.

\subsection{HDF5 Output for Disconnected Pieces}

The disconnected (quark loop and gluon) one-point functions are stored separately
by the inherited methods:

\begin{itemize}
  \item Quark 1pt: \texttt{save\_emt\_quark\_1pt\_hdf5} --
        saves per-noise-vector data (\texttt{raw/Tmunu\_pervec}, \texttt{raw/CHI\_pervec})
        and noise-averaged data (\texttt{avg/Tmunu/T11}--\texttt{T44}, \texttt{avg/CHI}).

  \item Gluon 1pt: \texttt{save\_emt\_gluon\_1pt\_hdf5} --
        saves momentum-projected gluonic EMT building blocks
        (\texttt{Tmunu/T11}--\texttt{T44}).
\end{itemize}

\subsection{File Organization}

The EMT output files are organized in subdirectories:
\begin{itemize}
  \item \texttt{EMTproton2pt/} -- proton two-point functions;
  \item \texttt{EMTproton3pt/} -- proton connected three-point functions;
  \item \texttt{EMTc/} -- stochastic quark one-point loops (shared with meson);
  \item \texttt{EMTg/} -- flowed gluon one-point functions (shared with meson).
\end{itemize}
The canonical EMTc loop name is source independent,
\texttt{<lat>.EMTc.<cfg>.<ama>.<sm>.h5}.  It stores all absolute insertion
times and is reused for every proton two-point source time on that
configuration.
Connected proton three-point file names append the sink kinematics as
\texttt{PX<px>PY<py>PZ<pz>dt<tsep>}.  A zero-momentum sink with one
source--sink separation at nine carries the suffix
\texttt{PX0PY0PZ0dt9.h5}; separations six, nine, and twelve are written as
three files with suffixes \texttt{dt6.h5}, \texttt{dt9.h5}, and
\texttt{dt12.h5}.  Each file stores a single-value \texttt{t\_separations}
attribute and keeps a length-one source--sink-separation axis in
\texttt{C3\_chi} and \texttt{C3\_Tmunu}.  This prevents runs with different
final momenta or source--sink separations from overwriting one another and
follows the nucleon TMD output convention.

% ---------------------------------------------------------------------------
\section{Disconnected Analysis}
% ---------------------------------------------------------------------------

\subsection{Combined Proton EMT}

The full proton matrix element of the EMT receives contributions from both
connected and disconnected diagrams:
\begin{equation}
  \langle p'|\Tmunu|p\rangle
  = \langle p'|\Tmunu|p\rangle_{\text{conn}}
  + \langle p'|\Tmunu|p\rangle_{\text{disc}}.
\end{equation}
The connected contribution is computed by \texttt{ProtonQuarkEMT.connected\_3pt}
as described in Sections~\ref{sec:C3_def}--\ref{sec:C3_total}.

\subsection{Disconnected Quark Contribution}
\label{sec:proton_disc_full_volume_source}

The disconnected quark contribution is evaluated using stochastic estimators.
For each noise vector $\xi^{(r)}$, one solves $D\eta^{(r)} = \xi^{(r)}$ and
computes the one-point function (spatially Fourier-transformed loop):
\begin{equation}
  L_{q,\,\mu\nu}^{(r)}(\momtrans,t)
  = -\frac{1}{2}\sum_{\mathbf{x}}
    \Phi_{\momtrans}(\mathbf{x})\;
    \xi^{(r)\dagger}(x)\,
    \gamma_{\nu}\big[D_{+\mu}-D_{-\mu}\big]\,
    \eta^{(r)}(x).
\end{equation}
At finite flow time the code evolves both stochastic fields with the same
four-dimensional fermion-flow kernel:
\begin{equation}
  \xi_f=K(t_f)\xi,
  \qquad
  \eta_f=K(t_f)D^{-1}\xi.
\end{equation}
Writing \(P_\tau\) for the absolute insertion-time projector, the estimator and
its expectation are
\begin{align}
  \widehat L(\tau,t_f)
  &=\xi^\dagger K^\dagger P_\tau\Gamma K D^{-1}\xi,
  \\
  \mathbb{E}_\xi[\widehat L(\tau,t_f)]
  &=\Tr\!\left[P_\tau\Gamma K(t_f)D^{-1}K^\dagger(t_f)\right].
  \label{eq:proton_full_volume_flowed_loop}
\end{align}
Thus the physical loop remains a spatial trace at fixed \(\tau\), but the
initial stochastic vector must span all Euclidean times because \(K(t_f)\)
spreads in the temporal direction.  A single initial-time projector does not
reproduce this finite-flow trace.  Complete time dilution would be unbiased
after summing all projectors, but the current workflow instead uses
full-volume counter-based \(Z_4\) sources and retains every insertion time.
Spatial HP and isotropic 4D HP both preserve this full-volume support; only
their probing signs differ.  A detailed derivation is provided in
\code{docs/EMT\_disconnected\_1pt/EMT\_disconnected\_1pt.tex}.

The disconnected three-point function is then obtained from the correlation with
the proton two-point function:
\begin{equation}
  C_{3,\,\mu\nu}^{\text{disc},\,q}(\momtrans,\tau)
  = \Big\langle C_{2}(t_{\text{sep}})\,
           \frac{1}{N_{v}}\sum_{r=1}^{N_{v}} L_{q,\,\mu\nu}^{(r)}(\momtrans,\tau)
    \Big\rangle
    - \big\langle C_{2}(t_{\text{sep}})\big\rangle\,
      \Big\langle \frac{1}{N_{v}}\sum_{r=1}^{N_{v}} L_{q,\,\mu\nu}^{(r)}(\momtrans,\tau)
      \Big\rangle.
\end{equation}
The subtraction of the product of averages removes the vacuum expectation value
of the loop.

\subsection{Disconnected Gluon Contribution}

The flowed gluonic EMT building block is built from the clover field strength:
\begin{equation}
  F_{\mu\nu}(x) = -\frac{i}{8}\Big(
    Q_{\mu\nu}(x) - Q_{\mu\nu}^{\dagger}(x)
    - \frac{1}{N_{c}}\Tr\big[Q_{\mu\nu}(x) - Q_{\mu\nu}^{\dagger}(x)\big]
  \Big),
\end{equation}
where $Q_{\mu\nu}$ is the average of the four plaquettes in the $(\mu,\nu)$ plane
(clover combination).  The gluonic EMT one-point building block is
\begin{equation}
  L_{g,\,\mu\nu}(\momtrans,t)
  = \frac{2}{V_{3}}\sum_{\mathbf{x}}
    \Phi_{\momtrans}(\mathbf{x})
    \sum_{\rho\neq\mu,\nu}
    \Tr_{c}\!\big[F_{\mu\rho}(x)\,F_{\nu\rho}(x)\big].
\end{equation}
This quantity is measured on the flowed gauge field at each flow time.  Its
disconnected contribution to the proton matrix element is
\begin{equation}
  C_{3,\,\mu\nu}^{\text{disc},\,g}(\momtrans,\tau)
  = \Big\langle C_{2}(t_{\text{sep}})\,
           L_{g,\,\mu\nu}(\momtrans,\tau)
    \Big\rangle
    - \big\langle C_{2}(t_{\text{sep}})\big\rangle\,
      \big\langle L_{g,\,\mu\nu}(\momtrans,\tau)\big\rangle.
\end{equation}

\subsection{Ringed-Fermion Normalization}

The connected and disconnected quark EMT outputs in this repository are
\emph{unringed} flowed bilinears.  They do not contain the ringed-fermion
normalization factor.  The ringed factor should be computed once in the analysis
stage from the quark one-point kinetic expectation value, then applied
consistently to both connected and disconnected quark EMT observables at the
same flow time.

The conventional ringed variables of Makino--Suzuki are normalized by
\begin{equation}
  Z_{\chi}^{\mathrm{ring}}(t)
  =
  \frac{-2\,\dim(R)\,N_f}
       {(4\pi)^2 t^2\,
        \left\langle
        \bar{\chi}(t,x)\overleftrightarrow{\not{D}}\chi(t,x)
        \right\rangle},
\end{equation}
where \(\dim(R)=N_c=3\) for fundamental QCD and \(N_f\) is the flavor
normalization convention used in the kinetic expectation value.  Since a quark
EMT bilinear contains one \(\chi\) and one \(\bar\chi\), the bilinear receives
one factor \(Z_{\chi}^{\mathrm{ring}}(t)\), not its square.

At zero momentum, the diagonal quark EMT trace provides the kinetic building
block:
\begin{equation}
  \sum_{\mu} L_{q,\,\mu\mu}(0,t)
  = -\frac{1}{2}\sum_{\mathbf{x}}
    \big\langle
      \bar{\chi}(t,\mathbf{x})\,
      \overleftrightarrow{\not{D}}\,
      \chi(t,\mathbf{x})
    \big\rangle,
\end{equation}
up to the lattice derivative convention and the spatial-volume normalization
documented in the one-point file.  In the current HDF5 convention,
\begin{equation}
  K_{\mathrm{code}}(t)
  =
  -2\left\langle
  \sum_{\mu=1}^{4}
  L_{q,\,\mu\mu}^{\mathrm{avg}}(q=0,t,\tau)
  \right\rangle_{\tau,\mathrm{cfg}},
\end{equation}
where the diagonal sum is reconstructed from \code{avg/Tmunu/T11},
\code{T22}, \code{T33}, and \code{T44}.  The quark 1pt run now performs this
algebraic extraction from the same stochastic EMT tensor and writes
\code{raw/kinetic\_pervec}, source bookkeeping,
\code{avg/kinetic\_spacetime}, and \code{flow\_times} to
\code{FlowedQuarkRinged/<lat>.FlowedQuarkRinged.<cfg>.<ama>.<sm>.h5}.
This adds no inversion, flow update, derivative contraction, or gather.

The companion is intentionally kinetic-only and contains no
configuration-local ringed factors.  Average \code{avg/kinetic\_spacetime}
over gauge configurations first, then evaluate the nonlinear factor; averaging
the inverses formed on separate configurations is incorrect.  The unflowed
step \(t=0\) must not be used for ringed normalization because the formula
contains \(t^{-2}\).  The standalone ringed-normalization workflow remains
available for dedicated high-statistics, spin--color dilution, block-output,
and resume runs.
The production quark one-point wrapper uses base-oriented HP-interval shards;
an explicit streaming finalizer validates complete base coverage before
publishing the canonical loop and kinetic companion.

For isospin-degenerate light \(u/d\) measurements, the same light-quark ringed
factor should multiply the connected \(U\) and \(D\) quark EMT observables.  If a
future workflow uses different quark masses, the analysis must either compute
flavor-specific factors from flavor-specific one-point files or explicitly form
the chosen flavor-average kinetic normalization before applying it.

% ===========================================================================
\part{Appendices}
% ===========================================================================

\appendix

% ---------------------------------------------------------------------------
\section{Summary of Einsum Index Conventions}
% ---------------------------------------------------------------------------

For quick reference, the index labels used throughout the code and this note are:

\begin{center}
\begin{tabular}{cl}
\toprule
Index & Meaning \\
\midrule
\texttt{w} & Checkerboard/parity index \\
\texttt{t, z, y, x} & Local spacetime coordinates on the current MPI rank \\
\texttt{x\_cb} & Checkerboarded spatial coordinate (in even-odd preconditioned data) \\
$i,j,k,l,m,n$ & Spin indices (range 0--3) \\
$a,b,c,d,e,f$ & Color indices (range 0--2 for SU(3)) \\
$g$ & Sink gamma matrix index (range 0--15) \\
$p$ & Momentum index \\
$q$ & Momentum transfer index \\
\texttt{pol} & Polarization index \\
\bottomrule
\end{tabular}
\end{center}

\noindent
Propagator layout: \texttt{[w, t, z, y, x\_cb, spin\_sink, spin\_src, color\_sink, color\_src]}.

\noindent
Sequential data layout: \texttt{[w, t, z, y, x\_cb, spin\_a, spin\_b, color\_a, color\_b]}.

% ---------------------------------------------------------------------------
\section{Gamma Matrix Reference}
% ---------------------------------------------------------------------------

The 16 bilinear operators and their PyQUDA gamma identifiers:

\begin{center}
\begin{tabular}{clr}
\toprule
Name & Description & PyQUDA ID \\
\midrule
\texttt{"I"}  & Scalar (identity)      & \texttt{gamma(0)}  \\
\texttt{"X"}  & $\gamma_{1}$           & \texttt{gamma(1)}  \\
\texttt{"Y"}  & $\gamma_{2}$           & \texttt{gamma(2)}  \\
\texttt{"SXY"}& $\sigma_{12}=[\gamma_{1},\gamma_{2}]/2$ & \texttt{gamma(3)} \\
\texttt{"Z"}  & $\gamma_{3}$           & \texttt{gamma(4)}  \\
\texttt{"SXZ"}& $\sigma_{13}$          & \texttt{gamma(5)}  \\
\texttt{"SYZ"}& $\sigma_{23}$          & \texttt{gamma(6)}  \\
\texttt{"T5"} & $\gamma_{4}\gamma_{5}$ & \texttt{gamma(7)}  \\
\texttt{"T"}  & $\gamma_{4}$           & \texttt{gamma(8)}  \\
\texttt{"SXT"}& $\sigma_{14}$          & \texttt{gamma(9)}  \\
\texttt{"SYT"}& $\sigma_{24}$          & \texttt{gamma(10)} \\
\texttt{"Z5"} & $\gamma_{3}\gamma_{5}$ & \texttt{gamma(11)} \\
\texttt{"SZT"}& $\sigma_{34}$          & \texttt{gamma(12)} \\
\texttt{"Y5"} & $\gamma_{2}\gamma_{5}$ & \texttt{gamma(13)} \\
\texttt{"X5"} & $\gamma_{1}\gamma_{5}$ & \texttt{gamma(14)} \\
\texttt{"5"}  & $\gamma_{5}$           & \texttt{gamma(15)} \\
\bottomrule
\end{tabular}
\end{center}

\noindent
The Dirac gamma matrices used for EMT derivative insertions:
\begin{align}
  \gamma_{1} &\leftrightarrow \texttt{gamma(1)}  & (D\text{-}\gamma\text{-}ID = 1),\\
  \gamma_{2} &\leftrightarrow \texttt{gamma(2)}  & (D\text{-}\gamma\text{-}ID = 2),\\
  \gamma_{3} &\leftrightarrow \texttt{gamma(4)}  & (D\text{-}\gamma\text{-}ID = 4),\\
  \gamma_{4} &\leftrightarrow \texttt{gamma(8)}  & (D\text{-}\gamma\text{-}ID = 8).
\end{align}

% ---------------------------------------------------------------------------
\section{Complete Two-Point Contraction Einsum Map}
% ---------------------------------------------------------------------------

For readers who want to verify the code against the formulas, here is the
complete mapping of the two-point function contraction steps.

\medskip\noindent
\textbf{Term 1:}
\begin{alignat}{2}
  &\text{Math:}\quad
    \eps_{abc}\,\eps_{def}\,
    (\Cmat\Gint)_{ij}\,(\Cmat\Gint)_{kl}\,
    [{\Gsink}_{g}]_{mn}\,
    \fwprop^{ik}_{ad}\,
    \fwprop^{jl}_{be}\,
    \fwprop^{mn}_{cf}
    &&\longrightarrow\; [g,p,t] \\
  &\text{Code:}\quad
    \begin{aligned}
      &\texttt{t1\_s1 = einsum("abc, wtzyxikad -> wtzyxikbcd")} \\
      &\texttt{t1\_s2 = einsum("ij, wtzyxjlbe -> wtzyxilbe")} \\
      &\texttt{term1\_sink = einsum("wtzyxikbcd, wtzyxilbe -> wtzyxklcde")} \\
      &\texttt{term1\_p3 = einsum("gmn, wtzyxmncf -> gwtzyxcf")} \\
      &\texttt{t1\_f1 = einsum("def, wtzyxklcde -> wtzyxklcf")} \\
      &\texttt{t1\_f2 = einsum("kl, wtzyxklcf -> wtzyxcf")} \\
      &\texttt{t1\_f3 = einsum("wtzyxcf, gwtzyxcf -> gwtzyx")} \\
      &\texttt{term1 = einsum("pwtzyx, gwtzyx -> gpt")}
    \end{aligned}
\end{alignat}

\medskip\noindent
\textbf{Term 2:}
\begin{alignat}{2}
  &\text{Math:}\quad
    \eps_{abc}\,\eps_{def}\,
    (\Cmat\Gint)_{ij}\,(\Cmat\Gint)_{kl}\,
    [{\Gsink}_{g}]_{mn}\,
    \fwprop^{ik}_{ad}\,
    \fwprop^{jn}_{be}\,
    \fwprop^{ml}_{cf}
    &&\longrightarrow\; [g,p,t] \\
  &\text{Code:}\quad
    \begin{aligned}
      &\texttt{t2\_s1 = einsum("abc, wtzyxikad -> wtzyxikbcd")} \\
      &\texttt{t2\_s2 = einsum("ij, wtzyxjnbe -> wtzyxinbe")} \\
      &\texttt{term2\_sink = einsum("wtzyxikbcd, wtzyxinbe -> wtzyxkncde")} \\
      &\texttt{term2\_p3 = einsum("gmn, wtzyxmlcf -> gwtzyxnlcf")} \\
      &\texttt{t2\_f1 = einsum("def, wtzyxkncde -> wtzyxkncf")} \\
      &\texttt{t2\_f2 = einsum("wtzyxkncf, gwtzyxnlcf -> gwtzyxkl")} \\
      &\texttt{t2\_f3 = einsum("kl, gwtzyxkl -> gwtzyx")} \\
      &\texttt{term2 = einsum("pwtzyx, gwtzyx -> gpt")}
    \end{aligned}
\end{alignat}

\medskip\noindent
\textbf{Final result:}
\begin{equation}
  \texttt{corr = - term1 - term2}.
\end{equation}

% ---------------------------------------------------------------------------
\section{List of Parameters}
% ---------------------------------------------------------------------------

The \texttt{parameters} dictionary expected by \texttt{ProtonQuarkEMT} and its
parent classes:

\begin{center}
\begin{tabular}{ll}
\toprule
Key & Description \\
\midrule
\texttt{qext} & List of momentum transfer vectors $q = (q_{x},q_{y},q_{z},0)$ \\
\texttt{pf} & Final-state proton momentum $p_{f}$ \\
\texttt{p\_2pt} & List of momenta for the two-point function \\
\texttt{CG\_GaussSmear} & Boolean: enable gauge-covariant Gaussian smearing \\
\texttt{pos\_boost} & Boost vector for forward source/sink smearing \\
\texttt{neg\_boost} & Boost vector for backward source/sink smearing \\
\texttt{boost\_in} & Boost vector for forward propagator (proton source) \\
\texttt{boost\_out} & Boost vector for sequential propagator (proton sink) \\
\texttt{width} & Gaussian smearing width \\
\texttt{pol} & List of polarization labels (e.g., \texttt{["PpUnpol"]}) \\
\texttt{t\_insert} & Maximum time separation (for TMD compatibility) \\
\texttt{save\_propagators} & Boolean: save intermediate propagators \\
\texttt{flow\_type} & Flow action type: \texttt{"wilson"} or \texttt{"symanzik"} \\
\texttt{flow\_epsilon} & Flow step size $\varepsilon$ \\
\texttt{flow\_steps} & Number of flow steps $N$ \\
\bottomrule
\end{tabular}
\end{center}

% ---------------------------------------------------------------------------
\section{Sign Convention Summary}
% ---------------------------------------------------------------------------

For quick reference, here are all the sign conventions used in the proton EMT code:

\begin{enumerate}
  \item \textbf{Two-point function:}
        $C_{2} = -(C_{2}^{(1)} + C_{2}^{(2)})$.  The minus sign is required for
        a positive-definite spectral decomposition of the proton correlator.

  \item \textbf{Sequential source (up insertion):}
        $D_{\text{total}} = -(R_{1} + R_{2} + R_{3} + R_{4})$.  The minus sign
        is consistent with the two-point function sign.

  \item \textbf{Sequential source (down insertion):}
        $D_{\text{(down)}} = D_{2} - D_{1}$.  The relative minus sign between the
        two terms is the standard antisymmetrization of the single down quark in
        the baryon.

  \item \textbf{Quark EMT first term:}
        $+\frac{1}{2}$ from the decomposition $\overleftrightarrow{D} =
        \overrightarrow{D} - \overleftarrow{D}$.

  \item \textbf{Quark EMT second term:}
        $-\frac{1}{2}$, because the left derivative contributes with opposite sign
        in the bilinear form.

  \item \textbf{Disconnected quark loop:}
        $-\frac{1}{2}$ overall factor, matching the connected convention.
        The fermion loop has an additional minus sign from the closed fermion loop.

  \item \textbf{Gluon EMT building block:}
        Factor $+2/V_{3}$ from the original code convention.  The clover field
        strength includes an extra factor of $-i$ to match the Euclidean convention.
\end{enumerate}


% ---------------------------------------------------------------------------
\section{Validation and Regression Evidence}
% ---------------------------------------------------------------------------

This section records the implementation-level validation that has been performed
for the connected proton EMT measurement.  These tests are not a replacement for
physics production checks such as continuum normalization, renormalization,
operator mixing, disconnected contributions, or final GFF extraction.  They are
designed to catch common implementation errors in the connected sequential-source
path: gauge covariance violations, wrong left-derivative placement, phase
convention mistakes, missing up-quark Wick contractions, and accidental changes
to the raw-only sequential builder.

\subsection{Validated Test Setup}

The most complete validation suite was run on Aurora with the S8T32 test gauge
\begin{verbatim}
software_gradientflow/Pyquda_Measurement/test_gauge/
  S8T32_wilson_b6.cg.1e-08.0
\end{verbatim}
using the PyQUDA develop environment
\begin{verbatim}
software_gradientflow/activate-pyquda-develop.sh
\end{verbatim}
with \texttt{backend="dpnp"}, \texttt{backend\_target="sycl"}, two MPI ranks, and
\texttt{mpi\_geometry=1.1.1.2}.  Parallel HDF5 was verified with
\texttt{h5py.get\_config().mpi == True}.  Quark source/sink smearing was disabled
for the gauge-covariance tests so that the tests isolate the EMT/sequential
contractions rather than the non-gauge-covariant boosted-smearing path:
\begin{verbatim}
EMT_PROTON_DISABLE_SMEARING=1
EMT_PROTON_GAUSS_SMEAR=0
EMT_PROTON_MG_BLOCK=none
\end{verbatim}
HYP gauge preprocessing was kept enabled, matching the normal workflow.

The suite output is archived under
\begin{verbatim}
/lus/flare/projects/StructNGB/xgao/run/test_gauge/EMT_proton/
  validation_suite_20260603_141831
\end{verbatim}
with logs, HDF5 outputs, temporary diagnostic drivers, and JSON summaries.  A
companion Markdown summary is kept in
\texttt{docs/proton\_EMT/validation\_summary.md}.

\subsection{Raw-Only Sequential Builder Regression}

The proton EMT path uses a raw-only sequential builder:
\texttt{create\_bw\_seq\_raw\_pyquda}.  The finalized proton sequential object is
then produced from the raw sequential propagator by the same
$\gamma_{5}$-hermiticity finalization used by the original qTMD-style builder.
The regression test compared
\begin{equation}
  \texttt{create\_bw\_seq\_pyquda(...)}
  \quad\text{against}\quad
  \texttt{\_final\_seq\_from\_raw\_prop(create\_bw\_seq\_raw\_pyquda(...))}.
\end{equation}
The measured difference was
\begin{equation}
  \max |\Delta| = 0,\qquad
  \frac{\max|\Delta|}{\max|C|} = 0,
\end{equation}
confirming that the memory-saving raw-only path does not change the finalized
sequential object.

\subsection{Gauge Covariance of the Connected Three-Point Function}

Gauge covariance was tested by comparing observables from the original gauge field,
a constant global SU(3) transformed field, and a deterministic local SU(3)
transformed field.  The tested gauge-invariant outputs were the two-point function
\texttt{C2}, the scalar/proton sequential contraction \texttt{C3\_chi}, and the EMT
insertion \texttt{C3\_Tmunu}.

For \texttt{FLOW\_STEPS=0}, \texttt{QMAX=0}, and \texttt{T\_SEPS=2}, the local
transform relative differences were
\begin{align}
  \delta_{\rm rel}(C_{2}) &= 1.1443\times 10^{-10},\\
  \delta_{\rm rel}(C_{3,\chi}) &= 5.0029\times 10^{-11},\\
  \delta_{\rm rel}(C_{3,T_{\mu\nu}}) &= 9.2440\times 10^{-12}.
\end{align}
The EMT tensor was explicitly symmetric to roundoff in the stored output:
\begin{equation}
  \max |T_{\mu\nu} - T_{\nu\mu}| = 0.
\end{equation}
This is the key regression for the left-acting derivative fix: before applying the
left derivative to the raw sequential propagator, the point-source local gauge
transform test isolated an $\mathcal{O}(1)$ failure in \texttt{C3\_Tmunu}; after
the fix, the difference is at solver/reduction roundoff level.

\subsection{Gradient-Flow Covariance}

The same baseline/global/local gauge-transform test was repeated with
\texttt{FLOW\_STEPS=1}.  The expected output shapes were observed:
\begin{align}
  \texttt{C3\_chi} &: [2,1,1,2,1,32],\\
  \texttt{C3\_Tmunu} &: [2,1,1,2,1,4,4,32].
\end{align}
The flow step changed the correlators nontrivially:
\begin{align}
  \max|C_{3,\chi}^{(1)} - C_{3,\chi}^{(0)}|
    &= 3.7547\times 10^{-5},\\
  \max|C_{3,T_{\mu\nu}}^{(1)} - C_{3,T_{\mu\nu}}^{(0)}|
    &= 7.5175\times 10^{-5}.
\end{align}
Across the full flow-step axis, the local transform relative difference for
\texttt{C3\_Tmunu} remained
\begin{equation}
  \delta_{\rm rel}(C_{3,T_{\mu\nu}}) = 9.2440\times 10^{-12}.
\end{equation}

\subsection{Momentum Phase Checks}

The \texttt{QMAX=1} smoke test produced 27 unique momenta with temporal component
zero, covering the full cube $[-1,0,1]^{3}$.  The $q=0$ slice of the \texttt{QMAX=1}
output matched the independent \texttt{QMAX=0} baseline exactly:
\begin{equation}
  \max |\Delta C_{3,\chi}(q=0)| = 0,\qquad
  \max |\Delta C_{3,T_{\mu\nu}}(q=0)| = 0.
\end{equation}
Additional \texttt{MomentumPhase} checks verified
\begin{equation}
  \Phi_{-q}(x) = \Phi_{q}(x)^{*},
\end{equation}
the corresponding conjugation relation for a real deterministic test field, and
self-consistency of the shifted-delta phase ratio; all three checks had zero
measured max difference in the test.

\subsection{Up-Quark Insertion Term Sanity}

The up insertion source was decomposed into the four Wick-contraction terms
\texttt{R1}, \texttt{R2}, \texttt{R3}, and \texttt{R4}.  The optimized source in
the production code was checked against the explicit decomposition:
\begin{equation}
  D_{\rm up} = -(R_{1}+R_{2}+R_{3}+R_{4}),
\end{equation}
with zero measured source-level difference.  After time slicing, momentum
projection, identity smearing, and linear sequential inversion, the raw full U
sequential propagator agreed with the sum of the four term-by-term raw sequential
propagators with
\begin{equation}
  \max|\Delta| = 1.1375\times 10^{-15},\qquad
  \delta_{\rm rel} = 6.4103\times 10^{-12}.
\end{equation}
The U and D raw sequential objects were finite, nonzero, and not accidentally
identical:
\begin{equation}
  \delta_{\rm rel}(U_{\rm raw},D_{\rm raw}) = 0.4890.
\end{equation}
This confirms that the two identical-up Wick-contraction channels are present in
the implemented U insertion and that the optimized source respects the explicit
four-term algebra.

% ---------------------------------------------------------------------------

\end{document}
